\documentclass[12pt,a4paper]{article}
\usepackage[utf8]{inputenc}
\usepackage[vietnamese]{babel}
\usepackage[margin=2cm]{geometry}
\usepackage{tcolorbox}

\pagestyle{plain}

\begin{document}

% --- HEADER 2 CỘT NGHỊ ĐỊNH 30/2020/NĐ-CP ---
\noindent
\begin{minipage}[t]{0.45\textwidth}
    \begin{center}
        \textbf{BỘ Y TẾ}\\
        \textbf{\underline{PHÂN VIỆN PHÁP Y QUỐC GIA}}\\[0.2cm]
        \small{Số: 99/PYBS-2026}\\
        \textit{V/v giám định bổ sung nguyên nhân tử vong}
    \end{center}
\end{minipage}
\hfill
\begin{minipage}[t]{0.52\textwidth}
    \begin{center}
        \textbf{CỘNG HÒA XÃ HỘI CHỦ NGHĨA VIỆT NAM}\\
        \textbf{\underline{Độc lập - Tự do - Hạnh phúc}}\\[0.2cm]
        \textit{Hà Nội, ngày 26 tháng 07 năm 2026}
    \end{center}
\end{minipage}

\vspace{0.8cm}

\begin{center}
    \textbf{\LARGE BÁO CÁO GIÁM ĐỊNH PHÁP Y BỔ SUNG}\\
    \textbf{\large (CHUYÊN ÁN TEST-99)}
\end{center}

\vspace{0.3cm}

\section*{I. PHÂN TÍCH DIỄN BIẾN THƯƠNG TỔN}
\begin{enumerate}
    \item \textbf{Tổn thương chấn thương sọ não nhẹ (vùng gáy):} Cú xô ngã của Tùng chỉ gây chấn thương cơ học tụ máu dưới da và \textbf{khiến nạn nhân bất tỉnh tạm thời trong 30 - 45 phút}, hoàn toàn không đe dọa tính mạng.
    \item \textbf{Tổn thương tử vong (vết đâm cổ):} Vết đâm đứt động mạch cảnh bởi mảnh thủy tinh vỡ xảy ra khi nạn nhân đang trong trạng thái nằm ngửa bất động trên sàn.
\end{enumerate}

\section*{II. KHUNG GIỜ TỬ VONG CHÍNH XÁC}
\begin{tcolorbox}[colback=red!10!white,colframe=red!75!black,title=\textbf{KẾT LUẬN KHUNG GIỜ TỬ VONG CHÍNH XÁC}]
\begin{center}
    Nạn nhân Nguyễn Văn Khang tử vong vì mất máu cấp chính xác vào khoảng:\\
    \textbf{\Large 21:00 đến 21:15 ngày 24/07/2026}
\end{center}
\end{tcolorbox}

\section*{III. SO KHỚP LOGIC LẠT TẨY (ALIBI CLASH PROOF)}

\begin{table}[h]
\centering
\begin{tabular}{|p{4.5cm}|p{5cm}|p{5cm}|}
\hline
\textbf{Đối tượng} & \textbf{Khung giờ hiện diện / Khai báo} & \textbf{Đánh giá Pháp y \& Logic} \\ \hline
\textbf{Tùng} (Bạn cũ) & Rời hiện trường lúc \textbf{20:15} (Camera cây xăng xác nhận). & Không thể là kẻ đâm chết Khang lúc \textbf{21:00}. \\ \hline
\textbf{Trần Thị Hà} (Bạn gái cũ) & Khai ở nhà cả tối nhưng lỡ lời mô tả bình trà vỡ. & \textbf{Có mặt sau 20:15 và đâm chết Khang lúc 21:00!} \\ \hline
\end{tabular}
\end{table}

\vspace{0.5cm}

\begin{tcolorbox}[colback=yellow!10!white,colframe=yellow!75!black,title=\textbf{KẾT LUẬN CHỦ CHỐT CHUYÊN ÁN}]
Hà bám theo nạn nhân, lén vào nhà sau khi Tùng rời đi, thấy Khang bất tỉnh + đọc được tin nhắn nhân tình mới $\rightarrow$ Cầm mảnh thủy tinh đâm chết Khang lúc 21:00.
\end{tcolorbox}

\vspace{0.8cm}

\noindent
\begin{minipage}[t]{0.45\textwidth}
    \textbf{\textit{Nơi nhận:}}\\
    \small{- Thủ trưởng Cơ quan CSĐT;}\\
    \small{- Viện KSMND Thành phố;}\\
    \small{- Ban Chuyên án TEST-99;}\\
    \small{- Lưu: VT, Phân viện PY.}
\end{minipage}
\hfill
\begin{minipage}[t]{0.5\textwidth}
    \begin{center}
        \textbf{CHỦ TỊCH HỘI ĐỒNG GIÁM ĐỊNH PHÁP Y}\\
        \textit{(Ký, ghi rõ họ tên và đóng dấu)}\\[1.5cm]
        \textbf{PGS. TS. Trần Mạnh Cường}
    \end{center}
\end{minipage}

\end{document}
